\documentclass[11pt]{article}
% Shared preamble for all daily exercises
% Update this file to change settings (padding, parindent, etc.) for all exercises

\usepackage[margin=0.75in]{geometry}
\usepackage{amsmath}
\usepackage{amssymb}

% Custom settings
\setlength{\parindent}{0pt}
\setlength{\parskip}{0.2cm}

\title{Daily Exercise}
\author{Dhruman Gupta}
\date{22nd April}

\begin{document}

% Common header for daily exercises
% This file can be included in other LaTeX documents

\maketitle

\noindent
\textbf{Course:} CS-3330, Theory of Computation

\vspace{0.2cm}

\noindent
\textbf{Question:}\noindent 
Prove that if $\Sigma_k^P = \Pi_k^P$, then for all $j > k$,

\[
\Sigma_k^P = \Sigma_j^P = \Pi_j^P
\]

\textbf{Answer:}

Assume that

\[
\Sigma_k^P = \Pi_k^P.
\]

We will show that the polynomial hierarchy collapses at level $k$.

First, recall that

\[
\Sigma_{k+1}^P = NP^{\Sigma_k^P}, \qquad \Pi_{k+1}^P = coNP^{\Sigma_k^P}.
\]

Now take any language $L \in \Sigma_{k+1}^P$. Then $L$ is decided by an NP machine with oracle in $\Sigma_k^P$.

As in the proof that $NP^{NP} = \Sigma_2^P$, an NP computation with a $\Sigma_k^P$ oracle gives one outer existential block, and then the oracle questions contribute a $\Pi_k^P$ condition. So $L$ has the form

\[
x \in L \iff \exists y \; P(x,y)
\]

where $P \in \Pi_k^P$.

But by assumption $\Pi_k^P = \Sigma_k^P$, so this is exactly a $\Sigma_k^P$ definition. Hence

\[
\Sigma_{k+1}^P \subseteq \Sigma_k^P.
\]

On the other hand, we always know that

\[
\Sigma_k^P \subseteq \Sigma_{k+1}^P.
\]

Therefore,

\[
\Sigma_{k+1}^P = \Sigma_k^P.
\]

Now take complements. Since $\Sigma_k^P = \Pi_k^P$, the class is closed under complement, and so

\[
\Pi_{k+1}^P = co\Sigma_{k+1}^P = co\Sigma_k^P = \Pi_k^P = \Sigma_k^P.
\]

Hence,

\[
\Sigma_{k+1}^P = \Pi_{k+1}^P = \Sigma_k^P.
\]

Now we prove the general case by induction on $j > k$.

The base case $j = k+1$ is already done.

Now assume that for some $j > k$,

\[
\Sigma_j^P = \Pi_j^P = \Sigma_k^P.
\]

Then

\[
\Sigma_{j+1}^P = NP^{\Sigma_j^P} = NP^{\Sigma_k^P} = \Sigma_{k+1}^P = \Sigma_k^P.
\]

Taking complements again,

\[
\Pi_{j+1}^P = co\Sigma_{j+1}^P = co\Sigma_k^P = \Pi_k^P = \Sigma_k^P.
\]

So

\[
\Sigma_{j+1}^P = \Pi_{j+1}^P = \Sigma_k^P.
\]

This completes the induction.

Therefore, for every $j > k$,

\[
\Sigma_k^P = \Sigma_j^P = \Pi_j^P.
\]

Hence proved.

\end{document}
